\hypertarget{appendix-a7fbaf6d1b1c}{}\label{referencestart:appendix-a7fbaf6d1b1c}
\section{Forward and backward error}\label{appendix-a7fbaf6d1b1c-forward-and-backward-error}

\textbf{Curriculum depth:} Introduced · \textbf{First seen:} C005 · \textbf{Latest development:} C005

\subsection{Definition and mechanism}\label{appendix-a7fbaf6d1b1c-definition-and-mechanism}

Forward error measures output discrepancy. Backward error measures how much the input must change for the computed output to be exact. Stability concerns algorithm errors; conditioning concerns problem sensitivity. Small backward error can coexist with large forward error.

\subsection{Prerequisites}\label{appendix-a7fbaf6d1b1c-prerequisites}

\begin{itemize}
\tightlist
\item
  \hyperref[appendix-3ae74e89b4e6]{Absolute and relative error} --- Discrepancy, magnitude and reference units define sensitivity and threshold significance.
\item
  \hyperref[appendix-8f6a01f284a7]{Conditioning and singularity} --- Input sensitivity is distinct from errors an implementation introduces.
\end{itemize}

\subsection{Used by}\label{appendix-a7fbaf6d1b1c-used-by}

Connect later dependent concepts here when their lessons arrive. - \hyperref[appendix-f8657bac66e5]{Matrix norms and numerical error} --- Extends this foundation into linear-systems and decomposition analysis. - \hyperref[appendix-d83b0ad1be7f]{A posteriori error bounds for contractive iteration} - \hyperref[appendix-aa2bf0d79639]{Algorithmic stability} - \hyperref[appendix-31b8c295e05f]{Numerical error budgets} \#\# Recall and next depth

Reconstruct the definition's example and name the assumptions that make it valid.

\textbf{Pending:} Linear algebra will develop concrete residuals, norms and backward-stability proofs.

\subsection{Source and navigation}\label{appendix-a7fbaf6d1b1c-source-and-navigation}

\hyperref[01-foundations-i-ch001-when-arithmetic-meets-uncertainty]{M005 teaching chapter} · \hyperref[appendix-fd090d982a15-m005---parts-77-88]{Part definitions} · \hyperref[appendix-4ad47bdd269a]{Dependency Map}
